\section{Cronograma y entregables}

\subsection{Plan de trabajo}

El estudio cualitativo se desarrolla en 10 semanas, del 19 de octubre al 24 de diciembre de 2026, y se entrega la última semana de diciembre.

\begin{table}[ht]
\centering
\small
\caption{Carta Gantt sintética del proyecto}
\label{tab:gantt}
\begin{tabular}{p{5.2cm}cccccc}
\hline
\textbf{Fase / etapa} & \textbf{\shortstack{Sem 1--2\\19--30 oct}} & \textbf{\shortstack{Sem 3--7\\2 nov--4 dic}} & \textbf{\shortstack{Sem 8--9\\7--18 dic}} & \textbf{\shortstack{Sem 10\\21--24 dic}} & \textbf{2027} \\
\hline
Fase 1: \textit{Desk research} con el catastro de JxR y co-diseño & $\blacksquare$ & & & & \\
Pauta validada con JxR y trámite de ética & $\blacksquare$ & & & & \\
Fase 2: Terreno con empresas, mujeres y ecosistema & & $\blacksquare$ & & & \\
Fase 3: Análisis QCA y comparación de casos & & & $\blacksquare$ & & \\
Informe final y entrega del cualitativo & & & & $\blacksquare$ & \\
Fase 4: Taller con JxR y aplicación del kit cuantitativo (según recursos) & & & & & $\blacksquare$ \\
\hline
\end{tabular}
\end{table}

\subsection{Paquete de entregables para la Directiva}

JxR recibirá cinco productos integrados:

\begin{enumerate}
    \item \textbf{Informe técnico y ejecutivo final (cualitativo, última semana de diciembre de 2026):} resultados del QCA, la reconstrucción de procesos de permanencia, los hallazgos de las viñetas pareadas por género y las conclusiones del estudio.
    \item \textbf{Matriz de tipología de empresas y guía B2B de abordaje y vinculación:} perfiles empresariales por disposición de adopción y manual práctico con argumentos para superar puntos de veto organizacionales.
    \item \textbf{Evaluación y propuesta de ajuste a los modelos de JxR:} diagnóstico de los ajustes requeridos por tipo de empresa para Bolsa Laboral, OTEC y CET.
    \item \textbf{Plan de rediseño operativo y hoja de ruta a 90 días:} recomendaciones operativas para mejorar la vinculación, la adopción y la permanencia laboral.
    \item \textbf{Kit metodológico cuantitativo:} cuestionario online, diseño muestral de unos 200 casos con cuotas por macrozona, plan de análisis y ejercicio comparativo hombre--mujer. Se presenta en 2027 y su aplicación depende de los recursos disponibles.
\end{enumerate}
